\maketitle

\begin{abstract}
Retrieval-augmented generation (RAG) can expose sensitive records before a language model produces any visible answer. We study where authorization should be enforced in a healthcare RAG pipeline and separate two outcomes that are often conflated: unauthorized context exposure during retrieval and unauthorized disclosure in generated text. Using deterministic Synthea v4.0.0 data, we evaluate 3,000 security cases (9,000 architecture executions) across prompt-only, pre-retrieval access control, and a risk-aware variant. Pre-retrieval authorization reduces unauthorized context exposure from 1,750/1,750 to 0/1,750, while the prompt-only canary disclosure rate is only 5/1,750 and its 5-to-0 reduction is not conventionally significant (exact McNemar $p=0.0625$). A controlled $k=1$ follow-up resolves an apparent utility penalty: all three architectures obtain 895/1,000 authorized-task successes when retrieval composition is held identical, showing that the earlier gap was caused by distractor composition rather than the authorization gate in this benchmark. A label-independent risk-feature redesign further shows that structural ACL denial generalizes to held-out phrasing, whereas heuristic suspicion triage does not: challenge-or-deny rates fall from 88.0\% to 16.7\% for fuzzy control and from 100\% to 16.0\% for deterministic rules. Finally, a nested 1K--100K retrieval study shows strong growth in unfiltered and proportional-scope retrieval cost and distractor similarity, while fixed-scope conditions serve as construction controls. The results support authorization-before-retrieval as a structural security boundary, but do not support claims that fuzzy triage outperforms transparent rules or that zero observed leakage proves zero real-world risk.
\end{abstract}

\begin{IEEEkeywords}
retrieval-augmented generation, access control, healthcare security, authorization, prompt injection, synthetic health records, formal verification
\end{IEEEkeywords}
